% !TEX root = ../matan.tex

\section{Теорема об обратном отображении}

\subsection*{Формулировка}

\begin{Theorem}[required]{Об обратном отображении}{}
	Пусть $f\colon\Omega\to\RR^d$, $\Omega\subset\RR^d$ открыто, $x\in\Omega$,
	и $f$ непрерывно дифференцируемо в окрестности точки $x$ (матрица-функция
	$f'(\xi)$ непрерывна при $|\xi-x|<r$). Если матрица Якоби $f'(x)$ обратима, то
	найдутся окрестность $U_x\ni x$ и окрестность $V\ni f(x)$ такие, что:
	\begin{itemize}
		\item $f$ --- биекция $U_x\to V$;
		\item обратное отображение $g=f^{-1}\colon V\to U_x$ непрерывно дифференцируемо;
		\item для каждой точки $\xi\in U_x$ с $y=f(\xi)$ справедливо
		$g'(y)=\bigl(f'(\xi)\bigr)^{-1}$.
	\end{itemize}
\end{Theorem}

\begin{Remark}{Нормировка}{}
	Без ограничения общности считаем $x=0$, $f(0)=0$ и $f'(0)=I_d$ (единичная
	матрица). Общий случай сводится к этому переходом к отображению
	$\tilde f=(f'(0))^{-1}\circ f$, у которого $\tilde f'(0)=I_d$; обратимость и
	дифференцируемость при этом сохраняются.
\end{Remark}

\subsection*{Вспомогательная лемма}

\begin{Lemma}{Возмущение тождественного отображения}{}
	Пусть $\vphi\colon B(0,r)\to\RR^d$ удовлетворяет условию Липшица с константой
	$0<c<1$: $\ \|\vphi(u)-\vphi(v)\|\le c\,\|u-v\|$. Положим $\Psi=\mathrm{id}+\vphi$,
	то есть $\Psi(u)=u+\vphi(u)$. Тогда
	\[
	(1-c)\,\|u-v\|\le\|\Psi(u)-\Psi(v)\|\qquad\forall u,v\in B(0,r),
	\]
	так что $\Psi$ инъективно. Если, кроме того, $\vphi(0)=0$, то
	\[
	B\bigl(0,(1-c)r\bigr)\subset\Psi\bigl(B(0,r)\bigr)\subset B\bigl(0,(1+c)r\bigr).
	\]
\end{Lemma}

\begin{proof}
	\emph{Нижняя оценка.} По неравенству треугольника
	\[
	\|u-v\|-\|\Psi(u)-\Psi(v)\|\le\|\vphi(u)-\vphi(v)\|\le c\,\|u-v\|,
	\]
	откуда $(1-c)\|u-v\|\le\|\Psi(u)-\Psi(v)\|$. При $u\ne v$ правая часть положительна,
	значит $\Psi(u)\ne\Psi(v)$ --- инъективность.

	\emph{Левое вложение.} Пусть $\vphi(0)=0$ и $\|w\|<(1-c)r$; найдём прообраз
	$\Psi(u)=w$ в $B(0,r)$. Уравнение $u+\vphi(u)=w$ равносильно
	$u=w-\vphi(u)$, то есть поиску неподвижной точки отображения $T(u)=w-\vphi(u)$.
	Выберем $r'<r$ так, что $\|w\|\le(1-c)r'$. На замкнутом шаре $\overline{B}(0,r')$
	отображение $T$ переводит шар в себя:
	\[
	\|T(u)\|\le\|w\|+\|\vphi(u)\|=\|w\|+\|\vphi(u)-\vphi(0)\|\le(1-c)r'+c r'=r',
	\]
	и является сжимающим: $\|T(u_1)-T(u_2)\|=\|\vphi(u_1)-\vphi(u_2)\|\le c\|u_1-u_2\|$.
	По теореме Банаха неподвижная точка существует --- это искомый прообраз.

	\emph{Правое вложение.} Для $\|u\|<r$:
	$\|\Psi(u)\|\le\|u\|+\|\vphi(u)-\vphi(0)\|\le(1+c)\|u\|<(1+c)r$.
\end{proof}

\subsection*{Доказательство теоремы}

\begin{proof}
	Считаем $f(0)=0$, $f'(0)=I_d$. Положим $\vphi(\xi):=f(\xi)-\xi$; тогда
	$\vphi'(\xi)=f'(\xi)-I_d$ и $\vphi'(0)=0$.

	\textbf{Шаг 1 ($\vphi$ --- сжатие).} Так как $f\in C^1$, матрица-функция $\vphi'$
	непрерывна, и из $\vphi'(0)=0$ следует, что найдётся $r'>0$, при котором
	$\|\vphi'(\zeta)\|\le\tfrac12$ для всех $\|\zeta\|\le r'$. По теореме о среднем для
	вектор-функций (для дифференцируемого $\vphi$ выполнено
	$\|\vphi(\xi)-\vphi(\eta)\|\le\sup_{\zeta\in[\xi,\eta]}\|\vphi'(\zeta)\|\cdot\|\xi-\eta\|$)
	получаем
	\[
	\|\vphi(\xi)-\vphi(\eta)\|\le\tfrac12\,\|\xi-\eta\|,\qquad \xi,\eta\in B(0,r').
	\]

	\textbf{Шаг 2 (биекция).} Применим лемму к $\Psi=\mathrm{id}+\vphi=f$ с $c=\tfrac12$:
	отображение $f$ инъективно на $B(0,r')$, и
	$B\bigl(0,\tfrac12 r'\bigr)\subset f\bigl(B(0,r')\bigr)$. Положим
	\[
	V:=B\bigl(0,\tfrac12 r'\bigr),\qquad U_x:=B(0,r')\cap f^{-1}(V).
	\]
	Тогда $f\colon U_x\to V$ --- биекция; обозначим $g:=f^{-1}\colon V\to U_x$.

	\textbf{Шаг 3 (дифференцируемость $g$).} Зафиксируем $y\in V$, $\xi=g(y)$ (так что
	$f(\xi)=y$); матрица $f'(\xi)$ обратима (при достаточно малом $r'$, по
	непрерывности $f'$ и обратимости $f'(0)=I_d$). Дадим $y$ приращение $\theta$,
	пусть $\tau$ --- соответствующее приращение прообраза:
	$g(y+\theta)=\xi+\tau$, $f(\xi+\tau)=y+\theta$, откуда $\theta=f(\xi+\tau)-f(\xi)$.
	Оценим
	\[
	\bigl\|g(y+\theta)-g(y)-(f'(\xi))^{-1}\theta\bigr\|
	=\bigl\|\tau-(f'(\xi))^{-1}\theta\bigr\|
	=\bigl\|(f'(\xi))^{-1}\bigl(f'(\xi)\tau-\theta\bigr)\bigr\|.
	\]
	Подставляя $\theta=f(\xi+\tau)-f(\xi)$ и пользуясь дифференцируемостью $f$:
	\[
	\le\|(f'(\xi))^{-1}\|\cdot\bigl\|f(\xi+\tau)-f(\xi)-f'(\xi)\tau\bigr\|
	=\|(f'(\xi))^{-1}\|\cdot o(\|\tau\|).
	\]
	Остаётся показать, что $o(\|\tau\|)=o(\|\theta\|)$. Из шага~1 и обратного
	неравенства треугольника
	\[
	\|\theta\|=\|f(\xi+\tau)-f(\xi)\|\ge\|\tau\|-\|\vphi(\xi+\tau)-\vphi(\xi)\|
	\ge\|\tau\|-\tfrac12\|\tau\|=\tfrac12\|\tau\|,
	\]
	так что $\|\tau\|\le 2\|\theta\|$ и $\theta\to0\iff\tau\to0$. Поэтому
	$o(\|\tau\|)=o(\|\theta\|)$, и
	\[
	\bigl\|g(y+\theta)-g(y)-(f'(\xi))^{-1}\theta\bigr\|=o(\|\theta\|),
	\]
	то есть $g$ дифференцируема в точке $y$ и $g'(y)=(f'(\xi))^{-1}$.

	\textbf{Шаг 4 ($g\in C^1$).} Формула $g'(y)=(f'(g(y)))^{-1}$ выражает $g'$ как
	композицию непрерывных отображений: $g$ непрерывно (как дифференцируемое),
	$\xi\mapsto f'(\xi)$ непрерывно по условию, а обращение матрицы --- непрерывная
	операция на множестве обратимых матриц. Значит, $g'$ непрерывна, $g\in C^1$.
\end{proof}
